% 几何约束
\section{几何约束：从"整个晶圆有多少bump"到"每条grid边能走多少线"}

第二类约束来自物理实现：链路不是抽象的——每条链路需要物理lane和信号bump来传输数据。
物理资源是有限的。几何约束给出$L_e$的\textbf{上界}。

\subsection{最粗糙的模型：全局lane预算}

将整块晶圆视为一个资源池：所有链路的总lane需求不能超过总bump数。
一条不等式，O(1)计算。但它忽略了一个关键事实：\textbf{bump按裸片分配，裸片之间不能借用。}
裸片$A$信号bump有余、裸片$B$不足，资源无法调剂。全局预算几乎不会成为绑定约束，但作为绝对上限仍有淘汰价值。

\subsection{更精细的模型：裸片级bump竞争——信号和电源的零和博弈}

每个裸片通过面阵列$\mu$bump与硅interposer连接。总bump数由裸片面积决定：$N_v^{\text{total}} = \eta A_v / p^2$。

核心矛盾：\textbf{信号bump和电源bump竞争同一物理面积。}
电源系统需要$P_v / (V_{\text{dd}} \cdot I_{\text{bump}})$个bump向裸片供电。
剩余的才分配给信号：

\begin{equation}\label{eq:bump-budget}
    N_v^{\text{sig}} = \frac{\eta A_v}{p^2} - \frac{P_v}{V_{\text{dd}} I_{\text{bump}}}
\end{equation}

裸片$v$上所有incident链路的lane需求受此预算约束：

\begin{equation}\label{eq:geom-per-die}
    \sum_{e \in \delta(v)} L_e \cdot \frac{B}{R_e} \le N_v^{\text{sig}} \quad \forall v \in \mathcal{V}
\end{equation}

其中$R_e$为链路$e$的每lane速率（取决于互联标准），$B/R_e$为需要的lane数。

写成矩阵形式：$\mathbf{M} \cdot \mathbf{L} \le \mathbf{b}$，$M_{v,e} = B/R_e$（当$e \in \delta(v)$）。

\textbf{注意一个精妙的耦合}：右手边$N_v^{\text{sig}}$本身包含$P_v = P_{0v} + \sum_{e \in \delta(v)} P_{\text{lane}} \cdot L_e \cdot B/R_e$。
$L_e$同时出现在不等式左边（lane需求）和右边（通过电源bump挤占信号bump）。
在实现中我们保守取$P_v$的额定最大值（TDP），使约束保持严格线性。
物理含义：\textbf{功耗越高，留给信号的bump越少——性能追求反噬了实现可行性。}

\subsection{更精细：布局感知布线容量}

裸片布局$p: \mathcal{V} \to \{0,\ldots,N-1\}^2$确定后，每条interposer grid边有走线容量上限。
链路$e$的Manhattan路径经过grid边$g$当且仅当$a_{g,e}=1$。
约束为：grid边$g$上的总走线需求$\sum_e a_{g,e} \cdot L_e \cdot B/R_e$不超过物理容量（布线层数$\times$每层密度）。

矩阵形式：$\mathbf{A}(p) \cdot \text{diag}(B/R_e) \cdot \mathbf{L} \le \mathbf{c}$。

\textbf{关键观察}：无论全局预算、per-die竞争还是布局感知布线，几何约束的形式都是$\mathbf{M} \cdot \mathbf{L} \le \mathbf{b}$。
精度升级只是矩阵$\mathbf{M}$的行数增加、数值更精确——\textbf{数学类型不变}。
